\chapter{QUADRATIC FORMS; ELEMENTARY GEOMETRY.}

The theory of quadratic forms, in its modern aspect, hardly goes back before the second half of the XVIIIth century, and, as we will see, it was developed mainly to respond to the needs of Arithmetic, Analysis and Mechanics. But the fundamental notions of this theory have in reality appeared from the beginnings of ``Euclidean'' geometry, of which they formed the frame. For this reason, its history can not be retraced without speaking, at least in a summary way, of the development of ``elementary geometry'' since antiquity. Of course, we will only be able to link up with the evolution of a few general ideas, and the reader must not expect to find here precise information on the history of such or such a particular theorem, about which it will suffice to refer back to specialised historical or didactic works. \(^{1}\) It goes without saying also, when we speak below of various possible interpretations of the same theorem in various algebraic or geometric languages, that we do not intend at all to say that these ``translations'' have at all times been as familiar as they are today; quite on the contrary, it is the principal gaol of this Note to show how, very gradually, mathematicians have become conscious of these relationships between questions often of very different aspect; we will also have to show how, doing this, they have been led to put some coherence in the mass of theorems from geometry bequeathed by the ancients, and finally to try to delimit exactly what should be understood by ``geometry''.

If one puts aside the discovery, by the Babylonians, of the formula for the solution of equations of the second degree ({[}232{]}, pp.~183-189) it is therefore under their geometric disguise that the birth of the principal concepts of the theory of quadratic forms must be noted. These appear first as squares of distances (in the plane or space of three dimensions) and the corresponding notion of ``orthogonality'' is introduced by means of the right angle, defined by Euclid as half the straight angle (Elements, Book I, Def. 10); the notions of distance and right angle being linked by the theorem of Pythagoras, the keystone of the Euclidean edifice. \(^{2}\) The idea of angle seems to have been introduced very early in Greek mathematics (which no doubt had received it from the Babylonians, broken in to the use of angles by their long astronomical experience). It is known that in the classical period, only angles less than 2 right angles are defined (the ``definition'' of Euclid is besides as vague and unusable as the ones he gives of straight line and plane); the notion of direction is not brought out, although Euclid uses (without axiom or definition) the fact that a straight line divides the plane into two regions, which he distinguishes carefully when it is necessary. \(^{3}\) At this stage, the idea of the group of plane rotations only comes to light in a very imperfect way, by addition (introduced, it also, without explanation by Euclid) of non directed angles of half-lines, which is only defined, in principle, when the sum is at most equal to two right angles. \(^{4}\) As for trigonometry, it is disdained by geometers and abandoned to surveyors and astronomers; it is these latter (Aristarchus, Hipparchus, Ptolemy especially {[}255{]}) who establish the fundamental relations between the sides and angles of a right angled triangle (plane or spherical) and draw up the first tables (they consist of tables giving the chord of the arc cut out by an angle \(\theta < \pi\) on a circle of radius r, in other words the number \(2r \sin \frac{\theta}{2}\) ; the introduction of the sine, more easily handled, is due to Hindu mathematicians of the Middle Ages); in the calculation of these tables, the formula for the addition of arcs, unknown at this time, is replaced by the equivalent use of the theorem of Ptolemy (going back perhaps to Hipparchus)

on quadrilaterals inscribed in a circle. It must be noted also that Euclid and Hero give propositions equivalent to the formula

\[
a ^ {2} = b ^ {2} + c ^ {2} - 2 b c \cos A
\]

between the sides and angles of an arbitrary plane triangle; but one can hardly see there the first appearance of the notion of a bilinear form associated with a metric form, for lack of the idea of a vector calculus which will only emerge in the XIXth century.

Displacements (or movements, the distinction between the two notions not being clear in antiquity --- nor even much later) are known to Euclid; but, for reasons of which we are ignorant, he seems to feel a clear repugnance to use it (for example in the ``case of equality of triangles'', where the impression is gained that he only uses the notion of displacement because he has not been able to formulate an appropriate axiom ({[}153 e{]}, v. I, pp.~225-227 and 249)); however, it is to the notion of displacement (rotation about an axis) that he has recourse for the definition of cones of revolution and of spheres (Elements, Book XI, def. 14 and 18), as well as Archimedes for that of quadrics of revolution. But the general idea of a transformation, applied to the whole of space, is more or less a stranger to mathematical thought before the end of the XVIIIth century;⁵ and before the XVIIth century, no trace is found either of the notion of composition of movements, nor a fortiori of the composition of displacements. That does not mean to say, of course, that the Greeks were not particularly aware of the ``regularities'' and ``symmetries'' of figures, which we now attach to the notion of groups of displacements: their theory of regular polygons, and even more that of regular polyhedra --- one of the most remarkable chapters of all their mathematics --- is there to prove the opposite.⁶

Finally, the last of the essential contributions of Greek mathematics, in the area which concerns us, is the theory of conics (as far as quadrics is concerned, the Greeks only know certain quadrics of revolution, and do not carry their study of them very far, except for the sphere). It is interesting to note here that, although the Greeks never had the idea of the fundamental principle of analytic geometry (essentially because of the lack of a manageable algebra), they commonly use, for the study of particular ``figures'', ``ordinates'' relative to two (or even more than two) axes in the plane (linked closely to the figure, which is one of the fundamental points where their method differs from that of Fermat and Descartes, whose axes are fixed independently of the figure under consideration). In particular, the first examples of conics (other than the circle) which come up in connection with the problem of the duplication of the cube, are the curves given by the equations \(y^{2}=ax, y=bx^{2}, xy=c\) (Menechme, a pupil of Eudoxus, middle of the IVth century); \(^{7}\) and it is the equation of conics (normally with reference to two oblique axes made up of a diameter and of the tangent to one of its points of contact with the curve) that is most often used in the study of problems relative to these curves (whereas the ``focal'' properties only play a very effacing role, contrary to what scholarly traditions going back only to the XIXth century would make us believe). From this vast theory, we must above all remember here the notion of conjugate diameters (already known by Archimedes), and the property that today serves as the definition of the polar of a point, given by Apollonius {[}153 b{]} when the point is exterior to the conic (the polar being therefore for him the straight line joining the points of contact of the tangents issuing from this point); from our point of view, these are two examples of ``orthogonality'' with respect to a quadratic form distinct from the metric form, but it is well understood that the link between these notions and the classical notion of perpendiculars could not possibly be conceived of at that time.

There is hardly any other progress to point out before Descartes and Fermat; but from the beginnings of analytic geometry, the algebraic theory of quadratic forms starts to emerge from its geometric crust: Fermat knows that an equation of the second degree in the plane represents a conic ({[}109{]}, v. I, pp.~100-102; French translation, v. III, pp.~84-101) and sketches analogous ideas about quadrics ({[}109{]}, v. I, pp.~111-117; French translation, v. III, pp.~102-108). With the development of analytic geometry in 2 and 3 dimensions during the course of the XVIIIth century two of the central problems of the theory appear (especially in connection with conics and quadrics): the reduction of a quadratic form to a sum of squares and the search for its ``axes'' with respect to the metric form. For conics, these two problems are too elementary to give rise to important algebraic progress; for an arbitrary number of variables, the first is solved by Lagrange in 1759, in connection with the maxima of functions of several variables ({[}191{]}, v. I, pp.~3-20). But this problem is almost immediately eclipsed by that of the search for axes, even before invariance of rank had been formulated; \(^{8}\) as for the law of inertia, it is only discovered around 1850 by Jacobi ({[}171{]}, v. III, pp.~593-598), who proves it by the same argument as nowadays, and Sylvester ({[}304{]}, v. I, pp.~378-381) who limits himself to stating it as quasi-obvious. \(^{9}\)

The problem of the reduction of a quadric to its axes already presents algebraic difficulties substantially greater than the analogous problem for conics; and Euler, who is the first to tackle it, is not in a position to prove the real character of the eigenvalues, which he admits after an attempt at justification without conviction ({[}108 a{]}, (1), v. IX, pp.~379-392). \(^{10}\) Given that this point is correctly established in about 1800 {[}140{]}, we must await Cauchy for the proof of the corresponding theorem for forms in an arbitrary number n of variables ({[}56 a{]}, (2), v. IX, pp.~174-195). It is also Cauchy who, at about the same time, proves that the characteristic equation giving the eigenvalues is invariant under all orthogonal changes of axes ({[}56 a{]}, (2), v. V, p.~252); \(^{11}\) but for n = 2 or n = 3, this invariance was intuitively ``obvious'' because of the geometric interpretation of eigenvalues by means of axes of the corresponding conic or quadric. Besides, during the course of research on this subject, the elementary symmetric functions of the eigenvalues have also occurred naturally (with various geometric interpretations, in connection notably with the theorems of Apollonius on conjugate diameters), and in particular the discriminant, which (known for a long time for n = 2 in connection with the theory of the equation of the second degree) appears for the first time for n = 3 with Euler ({[}108 a{]}, (1), v. IX, p.~382); this latter meets it in connection with the classification of quadrics (in expressing the condition for a quadric not to have a point at infinity) and does not mention its invariance with respect to orthogonal axis changes. But a bit later, with the beginnings of the arithmetical theory of quadratic forms with integral coefficients, Lagrange notes (for n = 2) a particular case of the invariance of the discriminant under linear but not orthogonal change of variables ({[}191{]}, v. III, p.~699), and Gauss establishes, for n = 3, the ``covariance'' of the discriminant for all linear transformations ({[}124 a{]}, v. I, pp.~301-302). \(^{12}\) Once the general formula for the multiplication of determinants was proved by Cauchy and Binet, the extension of the formula of Gauss to an arbitrary number of variables was immediate; it is that which, around 1845, will give the first impulsion to the general theory of invariants.

To the two notions that, amongst the Greeks, took the place of the theory of displacements --- that of movement and that of ``symmetry'' of a figure --- has just been added a third in the XVIIth and XVIIIth centuries with the problem of the changing of rectangular axes, which is substantially equivalent to this theory. Euler devotes several works to this question, attaching importance above all to obtaining parametric representations which are manageable for the formulae for changing axes. It is known what use Mechanics was going to make of the three angles which he introduces for this purpose for n = .3 ({[}108 a{]}, (1), v. IX, pp.~371-378). But he does not limit himself to that, considers in 1770 the general problem of orthogonal transformations for arbitrary n, remarks that the goal is reached here by introducing \(n(n - 1)/2\) angles as parameters, and finally, for n = 3 and n = 4, gives for rotations rational representations (as a function, respectively, of 4 homogeneous parameters and of 8 homogeneous parameters linked by a relation), which are none other than those obtained later by means of the theory of quaternions, and of which he does not indicate the origin ({[}108 a{]}, (1), v. VI, pp.~287-315). \(^{13}\)

On the other hand, Euler indicates also how to translate analytically the search for the ``symmetries'' of plane figures, and it is for this purpose that he is led to prove, substantially, that a plane displacement is a rotation, or a translation, or a translation followed by a symmetry ({[}108 a{]}, (1), v. IX, pp.~197-199). The rise of Mechanics at this time leads besides to the general study of displacements: but first of all it is only a question of ``infinitely small'' displacements tangent to continuous movements: they are apparently the only ones that occur in the researches of Torricelli, Roberval and Descartes on the composition of movements and the instantaneous centre of rotation for plane movements (cf.~p.~177). This latter is defined in a general way by Johann Bernoulli; d'Alembert in 1749, Euler the following year, extend this notion by proving the existence of an instantaneous axis of rotation for movements which leave a point fixed. The analogous theorem for finite displacements is not stated until 1775 by Euler {[}108 b{]}, in a memoir where he discovers at the same time that the determinant of a rotation is equal to 1; the following year, he proves the existence of a fixed point for plane similarities ({[}108 a{]}, (1), v. XXVI, pp.~276-285). But the work of Chasles must be awaited, starting in 1830 {[}60 a{]}, in order to obtain finally a coherent theory of finite and infinitely small displacements.

We arrive thus at what can be called the golden age of geometry, which is inserted grosso modo between the dates of the publication of the Géométrie descriptive of Monge (1795) {[}225{]} and of the ``Erlangen programme'' of F. Klein (1872) ({[}182{]}, v. I, pp.~460-497). The essential progress that we owe to this brusque renewal of geometry are the following:

\begin{enumerate}
\def\labelenumi{\Alph{enumi})}
\item
  The notion of the element at infinity (point, straight line or plane), introduced by Desargues in the XVIIth century {[}84{]}, but which is hardly expressed during the XVIIIth century except by abuse of language, is rehabilitated and systematically used by Poncelet {[}252{]} who thus makes of projective space the general framework of all geometrical phenomena.
\item
  At the same time, with Monge and especially Poncelet, the passage to complex projective geometry is effected. The notion of imaginary point, sporadically used during the XVIIIth century, is here exploited (concurrently with that of the point at infinity) to give statements independent of the ``case of the figure'' of real affine geometry. If at first the justifications brought in to support these innovations remained very awkward (especially those from the members of the school of ``synthetic'' geometry, where the use of coordinates ends up as being seen as a disgrace), what cannot be missed is recognising there, under the name of ``principle of contingent relations'' in Monge, or ``principle of continuity'' in Poncelet, the first germ of the idea of ``specialisation'' from modern algebraic geometry. \(^{14}\)
\end{enumerate}

One of the first results following on from these concepts is the remark that, in complex projective space, all the conics (respectively quadrics) which are non degenerate are of the same kind; which brings Poncelet to the discovery of ``isotropic'' elements:``Circles placed arbitrarily on a plane'', he says, ``are therefore not totally independent among themselves, as could be thought at first sight, they have ideally two imaginary points in common at infinity'' ({[}252{]}, v. I, p.~48). Later, he introduces in the same way the ``umbilical'', an imaginary conic at infinity common to all spheres ({[}252{]}, v. I, p.~370); and if he does not speak particularly of the isotropic generators of the sphere, at least he underlines explicitly the existence of rectilinear generators, real or imaginary, for all quadrics (ibid., p.~371); \(^{15}\) notions that his followers (notably

Plücker and Chasles), even more than he, make great use of, in particular in the study of ``focal'' properties of conics and quadrics.

\begin{enumerate}
\def\labelenumi{\Alph{enumi})}
\setcounter{enumi}{2}
\item
  The notions of point transformation and of composition of transformations are, they also, formulated in a general way and introduced systematically as means of proof. Apart from displacements and projections, only certain particular transformations were known up to that point: certain plane projective transformations of the type \(x' = a / x\) , \(y' = y / x\) , used by La Hire and Newton, the ``affinity'' \(x' = ax\) , \(y' = by\) of Clairaut and Euler ({[}108 a{]}, (1), v. IX, chap.~XVIII), and finally a few particular quadratic transformations, with Newton again, Maclaurin and Braikenridge. Monge, in his Géométrie descriptive, shows all the use that can be made of plane projections in 3 dimensional geometry. With Poncelet, one of the systematic procedures for proofs, used to surfeit, consists in reducing by projection properties of the conic to those of the circle (a method already used on occasion by Descartes and Pascal); and in order to go from a quadric to a sphere, he invents the first example of a projective transformation of space, the ``homology'' ({[}252{]}, v. I, p.~357); finally it is he also who introduces the first examples of birational transformations of a curve into itself. In 1827, Möbius ({[}223{]}, v. I, p.~217) (and independently Chasles in 1830 ({[}60 b{]}, p.~695)), define the most general linear projective transformations; at the same time appear inversion and other types of quadratic transformations, whose study will inaugurate the theory of birational transformations, which will develop in the second half of the XIXth century.
\item
  The notion of duality appears fully in the open and is found consciously linked to the theory of bilinear forms. The theory of poles and polars with respect to conics, which, since Apollonius, had only made some progress with Desargues and La Hire, is extended to quadrics by Monge, who, as well as his pupils, sees the possibility of transforming by this means known theorems into new results. \(^{16}\) But it is again due to Poncelet that the credit returns for having built up these remarks into a general method in the theory of transformations ``by reciprocal polars'', and of having made it into a particularly efficacious tool of discovery. A little later, notably with Gergonne, Plücker, Möbius and Chasles, the general notion of duality frees itself from the link with quadratic forms, still too close with Poncelet. In particular, Möbius, while examining the various possibilities for duality in space of 3 dimensions (defined by a bilinear form), discovers in 1833 duality with respect to an alternating bilinear form ({[}223{]}, v. I, pp.~489-515), \(^{17}\) especially studied, in the XIXth century, under the form of the theory of ``linear complexes'' and developed in relation to the ``geometry of straight lines'' and the ``Plücker co-ordinates'' introduced by Cayley, Grassmann and Plücker around 1860.
\item
  From the beginnings of projective geometry, the intensive study of the properties of classical geometry in their relationships with projective space had rapidly led to dividing them into ``projective properties'' and ``metric properties''; and it is no doubt not exaggerated to see in this separation one of the clearest manifestations, at this time, of what was to become the modern notion of structure. But Poncelet, who is the first to introduce this distinction and this terminology, is already conscious of what links these two types of properties; and, tackling in the Traité the problems concerning angles, whose properties ``seem not to be part of those that we have called projective \ldots, they follow on nonetheless in such a simple way'', he says, ``from the principles that form the basis {[}of this work{]} \ldots, that I do not believe that any other geometrical theory could lead to it in a manner which is both more direct and simpler. It will not be astonishing, if one considers that projective properties of figures are necessarily the most general among those that can belong to them; in such a way that they must encompass, as simple corollaries, all the other properties or particular relations of the area'' ({[}252{]}, v. I, p.~248). Truth to tell, after this declaration, it is a little surprising to see him tackle questions about angles in a very roundabout way, by linking them to focal properties of conics, instead of bringing in directly the cyclic points; and in fact, it is only 30 years later that Laguerre (still a pupil at the École Polytechnique) gave the expression for a right angle by means of the cross-ratio of these straight lines and isotropic straight lines with the same origin ({[}192{]}, v. II, p.~13). Finally, with Cayley ({[}58{]}, v. II, pp.~561-592) the fundamental idea that ``metric'' properties of a plane figure are none other than the ``projective'' properties of the figure extended by the cyclic points becomes clearly expressed --- a decisive way station towards the ``Erlangen programme''.
\item
  Non-Euclidean hyperbolic geometry, which will come to light around 1830, remains at first a little apart from the movement of which we are tracing the main features. Arising out of preoccupations of an essentially logical order touching on the foundations of classical geometry, this new geometry is presented by its inventors \(^{18}\) in the same axiomatic and ``synthetic'' form as Euclid's geometry, and without any link to projective geometry (whose introduction following the classical model seems even to be excluded a priori, since the notion of unique parallel disappears in this geometry); it is no doubt because of that that it hardly attracts, for a long time, any interest from the French, German or English schools of projective geometry. Also, when Cayley, in the fundamental memoir quoted earlier ({[}58{]}, v. II, pp.~561-
\end{enumerate}

\begin{enumerate}
\def\labelenumi{\arabic{enumi})}
\setcounter{enumi}{591}
\tightlist
\item
  has the idea of replacing cyclic points (considered as a ``tangentially degenerate'' conic) by an arbitrary conic (that he calls ``absolute''), he is not thinking at all of linking this idea with the geometry of Lobatschevsky-Bolyai, even though he indicates how his conception leads to new expressions for the ``distance'' between two points, and that he mentions its links with spherical geometry. The situation changes around 1870, when the non-Euclidean geometries, following the diffusion of the works of Lobatschevsky, and the publication of the works of Gauss and the inaugural lecture of Riemann, come to the fore of mathematical news. Following the path traced by Riemann, Beltrami, without knowing of the work of Cayley, finds again in 1868 expressions for distance as given by this latter, but in quite another context, by considering the interior of a circle as an image of a surface with constant curvature, in which the geodesics are represented by straight lines {[}18 a{]}; it is Klein, who two years later, makes (independently of Beltrami) a synthesis of these various points of view, which he completes by the discovery of elliptic non-Euclidean space ({[}182{]}, v. I, pp.~254-305). \(^{19}\)
\end{enumerate}

\begin{enumerate}
\def\labelenumi{\Alph{enumi})}
\setcounter{enumi}{6}
\tightlist
\item
  In the second half of the period which we are considering here, a period of critical reflection is inserted, during which the supporters of ``synthetic'' geometry, not content with having banished co-ordinates from their proofs, pretend to do without real numbers even in their axioms for geometry. The principal representative of this school is von Staudt, who succeeds essentially in bringing about this tour de force {[}325{]}, very admired in his time and even well into the XXth century; and if today the same importance is not attributed to ideas of this order, whose possibilities for fruitful application are revealed to be fairly thin, it must nonetheless be recognised that the efforts of von Staudt and his disciples contributed to clarifying the role of real or complex ``scalars'' in classical geometry, and to introducing in that very way the modern conception of geometries over an arbitrary base field.
\end{enumerate}

Around 1860, ``synthetic'' geometry was at its peak, but the end of its reign was approaching at great speed. Having remained heavy and ungracious during the whole of the XVIIIth century, analytic geometry, in the hands of Lamé, Bobillier, Cauchy, Plücker and Möbius, acquires at last the elegance and conciseness that will allow it to fight on equal terms with its rival. Especially, after about 1850, the ideas of group and invariant, formulated at last in a precise way, invade little by little the scene, and it is seen that the theorems of classical geometry are none other than the expression of identical relations between invariants or covariants of the group of similarities, \(^{20}\) in the same way that those of projective geometry express the identities (or ``syzygies'') between covariants of the projective group. It is this thesis that is magisterially stated by F. Klein in the famous ``Erlangen programme'' ({[}182{]}, v. I, pp.~460-497), where he advocates the abandonment of sterile controversies between the ``synthetic'' tendency and the ``analytical'' tendency; if, he says, the accusation levelled against this latter of giving a privileged role to an arbitrary system of axes ``was most often justified in what concerns the defective way in which the method of co-ordinates was used previously, it founders when it is a case of a rational application of this method . . . The domain of spatial intuition is not forbidden by the analytic method . . .'', and he underlines that ``one must not underestimate the advantage that a well adapted formalism brings to subsequent research, in that it goes in advance so to say of thought'' (loc. cit., pp.~488-490).

One concludes thus with a rational and ``structural'' classification of the theorems of ``geometry'' according to the group from which they arise: the linear group for projective geometry, the orthogonal group for metric questions, the symplectic group for ``linear complex'' geometry. But under this merciless clarity, classical geometry --- with the exception of algebraic geometry and differential geometry, \(^{21}\) henceforth constituted as autonomous sciences --- flares suddenly and loses all its brightness. Already the generalisation of methods based on the use of transformations had made the development of new theorems somewhat mechanical: ``Today,'', says Chasles in 1837 in his Aperçu historique, ``every one can put himself forward, take an arbitrary known truth, and submit it to various general transformation principles; he will draw out other truths, different or more general; and these will be susceptible to similar operations; so that one will be able to multiply, almost to infinity, the number of new truths deduced from the first \ldots{} So anyone can who wishes, in the present state of the science, generalise and create in Geometry; genius is no longer indispensable in order to add a stone to the structure'' ({[}60 b{]}, pp.~268-269). But the situation becomes much clearer with progress in the theory of invariants, which succeeds at last (at least for the ``classical'' groups) in formulating general methods allowing in principle the writing down of all the algebraic covariants and all their ``syzygies'' in a purely automatic way; a victory that, at the same time, signals the death, as a field of research, of the classical theory of invariants itself, and of ``elementary'' geometry, \(^{22}\) which has become in practice a simple dictionary for it. No doubt, nothing allows the foresight a priori, among the infinity of ``theorems'' that can thus be rolled out at will, as to which will be those whose statement, in the appropriate geometric language, will have a simplicity and an elegance comparable to those of the classical results, and there remains there a restricted domain where numerous amateurs continue to work with pleasure (the geometry of the triangle, the tetrahedron, of algebraic curves and surfaces of low degree, etc.). But for the professional mathematician, the mine is worked out, since there are no longer there any structural problems, susceptible of reverberating on other parts of mathematics; and this chapter of the theory of groups and invariants can be considered as closed until further notice. \(^{23}\)

Thus, after the Erlangen programme, Euclidean and noneuclidean geometries, from a purely algebraic point of view, became simple languages, more or less convenient, for expressing the results of the theory of bilinear forms, whose progress went in tandem with that of the theory of invariants. \(^{24}\) All that concerns the notion of rank of a bilinear form and the relationships between these forms and linear transformations is definitively clarified by the works of Frobenius ({[}119{]}, v. I, pp.~343-405). It is also due to Frobenius that we have the canonical expression of an alternating form on a free Z-module ({[}119{]}, v. I, pp.~482-544); however, the left symmetric determinant had already appeared with Pfaff, at the beginning of the century, in connection with the reduction of differential forms to a normal form; Jacobi, who, in 1827, takes up this problem again ({[}171{]}, v. IV, pp.~17-29), knows that a left symmetric determinant of odd order is null, and it is he who constructs the expression for the pfaffian and shows that it is a factor of the left symmetric determinant of even order; but he had not realised that this latter is the square of the pfaffian, and this point was only established by Cayley in 1849 ({[}58{]}, v. I, pp.~410-413). The notion of the symmetric bilinear form associated with a quadratic form is the most elementary case of the process of ``polarisation'', one of the fundamental tools of the theory of invariants. Under the name of ``scalar product'', this notion will know an immense prosperity, first with the vulgarisers of the ``vector calculus'', then, starting in the XXth century, thanks to the unsuspected generalisation that is brought to it by the theory of Hilbert spaces (see p.~212). It is also this latter theory that will bring to light the notion of the adjoint of an operator (which beforehand had hardly come out except in the theory of linear differential equations, and, in tensor calculus, by the waltz of the co- and contravariant indices under the baton of the metric tensor); it is that which finally will give all its character to the notion of hermitian form introduced first by Hermite in 1853 in connection with his arithmetic research ({[}159{]}, v. I, p.~237), but having remained on the edge of the great mathematical streams until around 1925 and the applications of complex Hilbert spaces to quantum theories.

The study of the orthogonal group and the group of similarities --- clearly conceived and treated as such since the middle of the XIXth century, and having become the heart of the theory of quadratic forms --- as well as that of the other ``classical'' groups (linear group, symplectic group and unitary group), takes up on the other hand a greater and greater importance. We can only mention here the essential role played by these groups, in the theory of Lie groups and differential geometry on the one hand, the arithmetic theory of quadratic forms (see for example {[}285{]} and {[}100{]}) on the other; \(^{25}\) it is to this circumstance, as well as to the extension of the concept of duality to the most diverse questions that is due the fact that there is hardly any modern mathematical theory where bilinear forms do not crop up in some form or other. We must in any case note that it is the study of rotation groups (in three dimensions) that led Hamilton to the discovery of quaternions {[}145 a{]}; this discovery is generalised by W. Clifford who, in 1876, introduces the algebras that bear his name, and proves that they are tensor products of quaternion algebras, or quaternion algebras and a quadratic extension ({[}65{]}, pp.~266-276). Found again four years later by Lipschitz {[}205 b{]}, who uses them to give a parametric representation of orthogonal transformations in n variables (generalising those that Cayley had obtained for n = 3 ({[}58{]}, v. I, pp.~123-126) and n = 4 (v. II, pp.~202-215) by the theory of quaternions), these algebras, and the notion of ``spinor'' that derives from them ({[}52 b{]} and {[}62 a{]}), would also know a great vogue in modern times by virtue of their use in quantum theories.

It remains for us finally to say a word about the evolution of ideas which led to the almost total abandonment of all restrictions on the ring of scalars in the theory of sesquilinear forms --- a common tendency in all modern algebra, but which has perhaps manifested itself here earlier than elsewhere. We have already pointed out the fruitful introduction of geometry over the field of complex numbers (which, besides, during the whole of the XIXth century, did not go without a perpetual and sometimes dangerous confusion between this geometry and real geometry); the clarity here comes especially from the axiomatic studies of the end of the XIXth century on the foundations of geometry {[}163 c{]}. In the course of this research, Hilbert and his emulators, notably, on examining the relationships between the various axioms, were led to construct appropriate counterexamples, where the ``base field'' (commutative or not) possessed more or less pathological properties, and they thus accustomed mathematicians to ``geometries'' of a quite new type. From an analytical point of view, Galois had already considered linear transformations where coefficients and variables took values in a finite prime field ({[}123{]}, p.~145); in developing these ideas, Jordan {[}174 a{]} is led in a natural way to consider the classical groups over these fields, groups whose intervention manifests itself in varied mathematical domains. Dickson, around 1900, extends the research of Jordan to all finite fields, and more recently, it has become apparent that a large part of the theory of Jordan-Dickson can be extended to an absolutely arbitrary ``base field''; this is due essentially to general properties of isotropic vectors and to Witt's theorem, which, trivial in the classical cases, were only established for an arbitrary base field in 1936 {[}337 a{]}. \(^{26}\)

But in thus pushing the study of sesquilinear forms towards an ``abstraction'' which is always greater, it has turned out to be extremely suggestive to keep as is the terminology which, in the case of space of 2 or 3 dimensions, came out of classical geometry, and to extend it to the n-dimensional case and even to infinite dimensional spaces. Passed over in its role as an autonomous and living science, classical geometry is thus transfigured into a universal language of contemporary mathematics, with a suppleness and a usefulness that are incomparable.

\chapter{TOPOLOGICAL SPACES.}

The notions of limit and continuity go back to antiquity; it would be impossible to make a complete history of them without studying systematically from this point of view, not only the mathematicians, but also the Greek philosophers and in particular Aristotle, nor also without following the evolution of these ideas across the mathematics of the Renaissance and the beginnings of the differential and integral Calculus. Such a study, which it would certainly be very interesting to undertake, would go far beyond the framework of this note.

It is Riemann who must be considered as the creator of topology, as of so many other branches of modern mathematics: it is in fact he who, first, sought to disengage the notion of topological space, conceived the idea of an autonomous theory of these spaces, defined the invariants (the ``Betti numbers'') which were to play the greatest role in the later development of topology, and gave its first applications to analysis (periods of abelian integrals). But the movement of ideas in the first half of the XIXth century had not happened without preparing the way for Riemann in more than one way. Indeed, the desire to settle mathematics on a firm base, which was the cause of so much important research during the whole of the XIXth century and up to our day, had led to the correct definition of the notion of a convergent series and of a sequence of numbers tending to a limit (Cauchy, Abel) and that of a continuous function (Bolzano, Cauchy). On the other hand, the geometric representation (by points in the plane) of complex numbers, or, as had been said until then, ``imaginary numbers'' (qualified sometimes also, in the XVIIIth century, as ``impossible'' numbers), a representation due to Argand and Gauss (see p.~161), had become familiar to the majority of mathematicians: it constituted progress of the same order as in our days the adoption of geometric language in the study of Hilbert space, and contained in embryo the possibility of a geometric representation of every object which was susceptible to continuous variation; Gauss, who in any case was naturally led to such conceptions by his research on the foundations of geometry, on non-Euclidean geometry, on curved surfaces, seems already to have had this possibility in view, for he uses the words ``size twice extended'' in defining (independently of Argand and the French mathematicians) the geometric representation of imaginary numbers ({[}124 a{]}, v. II, pp.~101-103 and pp.~175-178).

It is on the one hand his research on algebraic functions and their integrals, on the other hand his thoughts (largely inspired by his study of the works of Gauss) on the foundations of geometry, that led Riemann to formulate a programme of study which is exactly that of modern topology, and to give this programme a start towards its realisation. Here for example is how he expresses himself in his Theory of abelian functions ({[}259 a{]}, p.~91):

``In the study of functions that are obtained by the integration of exact differentials, some theorems of analysis situs are almost indispensable. Under this name, which was used by Leibniz, although perhaps with a somewhat different meaning, it is in order to include that part of the theory of continuous quantities which studies these quantities, not as independent of their position and measurable the ones by means of the others, but in abstracting every idea of measurement and studying solely their relationships of position and inclusion. I reserve the right to treat this object later, in a way that is completely independent of all measurement \ldots{}''.

And in his famous inaugural Lecture ``On the hypotheses that form the foundation of geometry'' ({[}259 a{]}, p.~272):

``\ldots The general notion of quantity several times extended, \(^{1}\) which contains that of spatial quantity as a particular case, has remained completely unexplored \ldots{} (p.~272)''.

``\ldots The notion of quantity presupposes an element susceptible of different determinations. Following on whether or not one can pass from one determination to another by continuos transitions, these determinations make up a continuous multiplicity (of which they will be called the points) or a discrete multiplicity (p.~273)''.

``\ldots The measurement consists of a superposition of quantities to compare; to measure, there must therefore be a means of bringing one quantity onto another. In the absence of such a means, it is not possible to compare two quantities except when one is part of the other\ldots{} The studies that can then be made about this subject constitute part of the theory of quantities, independent of the measurement, and where the quantities are not considered as having an existence independent of their position, nor as being expressible by means of a unit of measure, but as parts of a multiplicity. Such studies have become a necessity in many parts of mathematics, and in particular for the theory of analytic multiform functions \ldots{} (p.~274)''.

``\ldots The determination of the position in a given multiplicity is thus reduced, each time that it is possible, to a finite number of numerical determinations. There are, it is true, some multiplicities, in which the determination of the position requires, not a finite number, but an infinite sequence or a continuous multiplicity of determinations of quantities. Such multiplicities are made up for example by the possible determinations of a function in a given domain, the position of a figure in space, etc.'' (p.~276).

In this last sentence will be noticed the first idea of a study of functional spaces; besides, already in Riemann's Dissertation the same idea is to be found expressed ``The set of these functions'', he says in connection with the minimum problem known under the name of Dirichlet's principle, ``makes up a connected domain, closed in itself'' ({[}259 a{]}, p.~30), which is, in an imperfect form, nonetheless the embryo of the proof that Hilbert would give later of Dirichlet's principle, and even of the majority of the applications of functional spaces to the calculus of variations.

As we have said, Riemann gave a start to the carrying out of this grandiose programme, by defining the ``Betti numbers'', first of a surface ({[}259 a{]}, pp.~92-93), then (ibid., pp.~479-482; cf.~also {[}259 c{]}) of a multiplicity of an arbitrary number of dimensions, and in applying this definition to the theory of integrals; results which inaugurate algebraic Topology, a branch of mathematics whose development has not stopped accelerating since the beginning of the XXth century, and should not be traced here.

As for the general theory of topological spaces, such as was foreseen by Riemann, it was necessary in order for it to develop, that the theory of real numbers, of sets of numbers, of sets of points on the straight line, in the plane and in space, should first be studied more systematically than had been the case in Riemann's time: this study was linked, on the other hand, to researches (semi philosophical with Bolzano, essentially mathematical with Dedekind) on the nature of irrational numbers, as well as progress in the theory of functions of a real variable (to which Riemann himself brought an important contribution by his definition of the integral and his theory of trigonometric series, and which was the object, among others, of the work of du Bois-Reymond, Dini, Weierstrass); it was the work of the second half of the XIXth century and most particularly Cantor, who was the first to define (first on a straight line, then in Euclidean space of n dimensions) the notions of accumulation point, closed and perfect set, and who obtained the essential results on the structure of these sets on the straight line(cf.~p.~155): one should consult on this, not only the Works of Cantor {[}47{]}, but also his very interesting correspondence with Dedekind {[}48{]}, where will be found clearly expressed also the idea of the number of dimensions considered as a topological invariant.

Further progress of the theory is stated for example, in a semi-historical, semi-systematic form in the book of Schoenflies {[}275 a{]}: by far the most important was the Borel-Lebesgue theorem (proved first by Borel for a closed interval on the straight line and a countable family of open intervals covering it).

The ideas of Cantor had first met a fairly lively opposition (cf.~p.~28). At least his theory of sets of points on the straight line and in the plane was soon used and widely spread by the French and German schools of the theory of functions (Jordan, Poincaré, Klein, Mittag-Leffler, then Hadamard, Borel, Baire, Lebesgue, etc.): the first volumes of the Borel collection, in particular, each contain an elementary statement of this theory (see for example {[}32 a{]}). As these ideas were spread, people began to think in different ways about their possible application to sets, not of points anymore, but of curves or functions: an idea that saw the light of day, for example, already in 1884, in the title ``On the limit curves of a family of curves'' of a memoir of Ascoli {[}10{]}, and which expresses itself in a communication of Hadamard to the congress of mathematicians at Zurich in 1896 {[}141{]}; it is tightly linked also to the introduction of ``line functions'' by Volterra in 1887, and to the creation of the ``functional calculus'', or theory of functions where the argument is a function (on that, the work of Volterra on functional Analysis {[}322 b{]} can be consulted). On the other hand, in the famous memoir ({[}163 a{]}, v. III, pp.~10-37) where Hilbert, taking up again on this point the ideas of Riemann, proved the existence of the minimum in the Dirichlet principle and inaugurated the ``direct method'' of the calculus of variations, one sees clearly appearing the interest that there is in considering sets of functions where the Bolzano-Weierstrass principle is valid, that is to say every sequence contains a convergent sub sequence; such sets would soon indeed play an important role, not only in the calculus of variations, but in the theory of functions of a real variable (Ascoli, Arzelà) and in that of functions of a complex variable(Vitali, Carathéodory, Montel). Finally the study of functional equations, and most particularly the solution by Fredhom of the type of equation that bears his name {[}116{]}, made familiar the consideration of a function as an argument, and of a set of functions as the analogue of a set of points, in connection with which it was quite as natural to use geometric language as for the points of a Euclidean space of n dimensions (a space which, also, escapes ``intuition'', and for that reason, remained for a long time an object of mistrust for many mathematicians). In particular, the memorable works of Hilbert on integral equations {[}163 b{]} end with the definition and the geometric study of Hilbert space by Erhard Schmidt {[}274 b{]}, in complete analogy with Euclidean geometry (see p.~212).

Meanwhile the notion of axiomatic theory had taken an increasingly greater importance, thanks above all to numerous works on the foundations of geometry, among which those of Hilbert {[}163 c{]} exercised a particularly decisive influence; during this very same work, Hilbert had been led to state exactly, from 1902 ({[}163 c{]}, p.~180), a first axiomatic definition of the ``multiplicity twice extended'' in the sense of Riemann, a definition which constituted, he said, ``the foundation of a rigorous axiomatic treatment of analysis situs'', and already used neighbourhoods (in a meaning restricted by the exigencies of the problem to which Hilbert was then limiting himself).

The first attempts to bring out that which was in common to the properties of sets of points and of functions (without the intervention of a notion of ``distance''), were made by Fréchet {[}115 a{]} and F. Riesz {[}260 b{]}; but the first, starting from the notion of countable limit, does not succeed in constructing, for non metrizable spaces, an axiom system which is convenient and fruitful; at least he recognised the relationship between the Bolzano-Weierstrass principle and the Borel-Lebesgue theorem; it is in this connection that he introduced the word ``compact'', even though with a meaning somewhat different from that given to it today. As for F. Riesz, who started from the notion of point of accumulation (or rather, which comes to the same thing, of ``derived'' set), his theory was still incomplete, and remained besides as only a sketch.

With Hausdorff ({[}152 a{]}, chap.~7-8-9) general topology as it is understood today starts. Taking up again the notion of neighbourhood, he knew how to choose, among the axioms of Hilbert on neighbourhoods in the plane, those that could give to his theory at the same time all the precision and the generality that were desirable. The chapter where he develops its consequences remains a model of axiomatic theory, abstract but adapted to applications from the beginning. That was, quite naturally, the point of departure for subsequent research on general topology, and principally the work of the school of Moscow, directed for the greater part towards the problem of metrisation (cf.~p.~166): we must remember from this especially here the definition, by Alexandroff and Urysohn, of compact spaces (under the name of ``bicompact spaces''), then the proof by Tychonoff {[}313{]} of the compactness of products of compact spaces. Finally the introduction of filters by H. Cartan {[}53{]}, which while supplying a very valuable instrument in view of all kinds of applications (where it is an advantageous substitute for the notion of ``Moore-Smith convergence'' {[}227{]}), came, thanks to the theorem on ultrafilters, to achieve a clarification and simplification of the theory.

\chapter{UNIFORM SPACES.}

The principal notions and propositions relating to uniform spaces came out little by little from the theory of real variables, and were the object of systematic study only recently. Cauchy, looking to set up rigorously the theory of series (cf.~p.~153), took as departure point a principle that he seems to have considered as obvious, according to which a necessary and sufficient condition for the convergence of a sequence \((a_{n})\) is that \(|a_{n+p} - a_{n}|\) should be arbitrarily small whenever \(n\) is sufficiently large (see for example ({[}56 a{]}, (2), v. VII, p.~267)). With Bolzano {[}27 c{]}, he was without doubt one of the first to state explicitly this principle, and to recognise its importance: whence the name ``Cauchy sequence'' given to sequences of real numbers that satisfy the condition in question, and by extension to sequences \((x_{n})\) of points in a metric space such that the distance between \(x_{n+p}\) and \(x_{n}\) is arbitrarily small whenever \(n\) is sufficiently large; from there finally the name of ``Cauchy filter'' given to the generalisation of Cauchy sequences in uniform spaces.

When later on the intuitive notion of real number was no longer satisfactory, and that it was sought, in order to give Analysis a solid foundation, to define real numbers starting from rational numbers, it was precisely the Cauchy principle that furnished the most fruitful among the definitions suggested in the second half of the XIXth century; it is Cantor's definition ({[}47{]}, pp.~93-96) (developed also, among others, from the ideas of Cantor, by Heine {[}154 b{]}, and independently, by Méray), according to which one makes a real number correspond to every Cauchy sequence (``fundamental sequence'' in the terminology of Cantor) of rational numbers; the same real number will correspond to two Cauchy sequences of rational numbers \((a_{n})\) and \((b_{n})\) if \(|a_{n} - b_{n}|\) tends to zero, and in this case only. The essential idea here is that, from a certain point of view, the set \(\mathbf{Q}\) of rational numbers is ``incomplete'', and that the set of real numbers is the ``complete'' set that is deduced from \(\mathbf{Q}\) by ``completing'' it.

On the other hand, Heine, in work mainly inspired by the ideas of Weierstrass and Cantor, was the first to define uniform continuity for numerical functions of one or more real variables {[}154 a{]}, and proved that every numerical function, continuous on a closed bounded interval of R, is uniformly continuous there {[}154 b{]}: it is ``Heine's theorem''. This result is linked to the compactness of a closed bounded interval in R (``Borel-Lebesgue theorem'', cf.~p.~141), and the proof given by Heine of his theorem can also be used, with a few modifications, to prove the Borel-Lebesgue theorem (which seemed to several authors a sufficient reason for giving this theorem the name ``Heine-Borel theorem'').

The extension of these ideas to more general spaces was made when, first some special cases were studied, and then in general, metric spaces, where a distance (numerical function of pairs of points, satisfying certain axioms) is given and defines at the same time a topology and a uniform structure. Fréchet, who was the first to state the general definition of these spaces, recognised the importance of the Cauchy principle {[}115 a{]}, and introduced also for all metric spaces, the notion of pre compact (or ``totally bounded'' {[}115 a and b{]}) space. Hausdorff, who, in his ``Mengenlehre'' {[}152 a and b{]} developed a lot of the theory of metric spaces, recognised in particular that one can apply to these spaces the Cantor construction which was considered above, and so deduce, for every non ``complete'' metric space (that is to say where the Cauchy principle is not valid), a ``complete'' metric space.

Metric spaces are ``uniform spaces'' of a particular kind; uniform spaces have only been defined in a general way recently, by A. Weil {[}330 b{]}. Beforehand it was only known how to use the notions and results referring to ``uniform structure'' when metric spaces were being dealt with: which explains the important role played in a lot of modern work on topology by metric or metrizable spaces (and in particular by compact metrizable spaces) in questions where distance is not of any real use. Once the definition of uniform spaces is stated, there is no difficulty (especially when the notion of filter is also available) in extending to these spaces nearly all the theory of metric spaces, as it is stated for example by Hausdorff (and to extend in the same way, for example, to all compact spaces, the results set out for compact metric spaces in the Topology of Alexandroff-Hopf {[}4{]}). In particular, the completion theorem for uniform spaces is none other than the transposition, without any essential modification, of Cantor's construction for the real numbers.

\chapter{REAL NUMBERS.}

Every measure of size implies a diffuse notion of real numbers. From the point of view of mathematics, the origins of the theory of real numbers must go back to the progressive formation, in Babylonian science, of a system of numbering capable (in principle) of denoting values as near as required to every real number {[}232{]}. The possession of such a system, and the confidence in numerical calculations which can not but follow from this, terminated inevitably, in fact, with a ``naive'' notion of real numbers, which is hardly different from that which one finds again today (linked to the decimal system of numbering) in elementary teaching or with physicists and engineers; this notion does not allow itself to be defined with exactness, but it can be expressed by saying that a number is considered as defined by the possibility of obtaining approximate values and introducing these into the calculations: which, besides, implies necessarily a certain degree of confusion between the measures of size given by experience, which are naturally not susceptible to indefinite approximation, and such ``numbers'' as \(\sqrt{2}\) (supposing that there exists an algorithm for the indefinite approximation of this number).

A similar ``pragmatic'' point of view reappears therefore in all the schools of mathematics where calculating skills take over from the desire for rigour and theoretical preoccupations. It is these latter, on the contrary, which dominate Greek mathematics: so we owe to it the first rigorous and coherent theory of ratios of quantities, that is to say, essentially the real numbers; it is the end result of a series of discoveries on ratios and in particular on incommensurable ratios, of which it is difficult to exaggerate the importance in the history of Greek thought, but of which, in the absence of precise texts, we can only just discern the major lines. Greek mathematics in its beginnings is inseparably linked to speculations, partly scientific, partly philosophical and mystical, on proportions, similitiades and ratios, in particular the ``simple ratios'' (expressible by fractions with small numerator and denominator); and it was one of the characteristic tendencies of the Pythagorean school to claim to explain all by integers and the ratio of integers. But it was the Pythagorean school, precisely, that discovered the incommensurability of the side of a square with its diagonal (the irrationality of \(\sqrt{2}\) ): the first example, no doubt, of a proof of impossibility in mathematics; just the fact of stating this question implies a clear distinction between a ratio and its approximate values, and suffices to indicate the immense gap which separates Greek mathematicians from their predecessors. \(^{1}\)

We are badly informed about the movement of ideas that accompanied and followed this important discovery. \(^{2}\) We will limit ourselves to indicating summarily the principal ideas that are at the basis of the theory of ratios of quantities, a theory which, codified by the great mathematician Eudoxus (contemporary and friend of Plato), was definitively adopted by classical Greek mathematics, and is known to us through the Elements of Euclid {[}107{]}, where it is given a masterly exposition (in Book V of these Elements):

\begin{enumerate}
\def\labelenumi{\arabic{enumi})}
\item
  The word and the idea of number are strictly reserved for integers \textgreater{} 1 (1 is the monad and not a number properly speaking), to the exclusion, not only of our irrational numbers, but even of those that we call rational numbers, these being, for the Greek mathematicians of the classical era, ratios of numbers. There is there much more than a simple question of terminology, the word number being linked for the Greeks (and for the moderns until recent times) to the idea of a system with a double law of composition (addition and multiplication): the ratios of integers are conceived by the classical Greek mathematicians as operators, defined on the set of integers or on a part of this set (the ratio of p to q is the operator which, to N, makes correspond, if N is a multiple of q, the integer \(p.(N/q)\) ), and making up a multiplicative group, but not a system with a double law of composition. In this, the Greek mathematicians were separating themselves voluntarily from the ``logisticians'' or professional calculators, who had, like their Egyptian and Babylonian predecessors, no scruples at treating as numbers fractions or the sum of an integer and a fraction. It seems besides that they imposed on themselves this restriction on the idea of number for motives more philosophical than mathematical, and following on from the thoughts of the first Greek thinkers on the unit and the multiple, the unit not being capable of subdivision (in this system of thought) without losing by that selfsame action its character of unit. \(^{3}\)
\item
  The theory of quantities is axiomatically based, and at the same time for all kinds of quantities (one finds allusions to previous theories which, from what appears, treated separately lengths, areas, volumes, time, etc.). The quantities of one kind are characterised by the fact that they are capable of being compared (that is to say that a definition of equality, which is properly speaking an equivalence, and the relations \textgreater{} and \textless{} are assumed), of being added and subtracted \(A + B\) is defined, and A - B if A \textgreater{} B), and satisfy the axiom of ``Archimedes'' so-called; this is clearly conceived, from the beginning, as a keystone of the edifice (it is in fact indispensable for all axiomatic characterisations of real numbers); it is a pure accident that it has been given the name of Archimedes, and this latter insists, in the introduction to his ``Quadrature of the Parabola'' ({[}5 b{]}, v. II, p.~265), on the fact that this axiom has been used by his predecessors, that it plays an essential role in the work of Eudoxus, and that its consequences are no less sure than the determinations of areas and volumes made without its help. \(^{4}\)
\end{enumerate}

It is easy to see that, from this axiomatic foundation, follows necessarily the theory of real numbers. It must be noted that, for Eudoxus, the quantities of a given kind form a system with one internal law of composition (addition), but that this system has an external law of composition with as operators the ratios of quantities, these being conceived of as forming an abelian multiplicative group. A and \(A'\) being quantities of the same kind, and similarly B and \(B'\) , the ratios of A and \(A'\) and of B and \(B'\) are defined as equal if, whatever the integers m and \(m'\) , \(mA < m'A'\) forces \(mB < m'B'\) and \(mA > m'A'\) forces \(mB > m'B'\) ; inequalities between ratios are defined in an analogous manner. That these ratios form a domain of operators for every kind of quantity is equivalent to the axiom (not made explicit but often used in the drafting of Euclid) of the existence of a fourth proportional: a ratio \(A/A'\) being given and \(B'\) being given, there exists a B, of the same kind as \(B'\) , such that \(B/B' = A/A'\) . Thus the highly original idea of Eudoxus allowed identification between the domains of operators defined for all kinds of quantities;⁵ in an analogous manner, the set of integers (see above) can be identified with a subset of the set of ratios of quantities, namely with the set of rational ratios (ratios of commensurable quantities); however, because of the fact that these ratios, in as much as they are operators on the integers, are (in general) defined only for a subset of the set of integers, it was still necessary to develop their theory separately (Book VII of Euclid).

The domain of universal operators thus constructed was therefore for Greek mathematicians the equivalent of what is for us the set of real numbers; it is clear besides that with the addition of quantities and the multiplication of ratios of quantities, they possessed the equivalent of what for us is the field of real numbers, even though in a much less handy form. \(^{6}\) It is possible, on the other hand, to ask oneself if they had conceived of these sets (set of quantities of a given kind, or sets of ratios of quantities) as complete in our sense; it is not easy to see, otherwise, why they would have assumed (without even feeling the need to turn it into an axiom) the existence of a fourth proportional; furthermore, certain texts seem to refer to ideas of this kind; finally they certainly assume as obvious that a curve, capable of being described by a continuous movement, can not pass from one side to the other of a straight line without cutting it, a principle that they used for example in their research on the duplication of the cube (construction of \(\sqrt[3]{2}\) by the intersection of curves) and which is essentially equivalent to the property under consideration; however, the texts that we possess do not allow us to know with complete precision their ideas on this point.

Thus therefore is the state of the theory of real numbers in the classical era of Greek mathematics. As admirable as the construction of Eudoxus was, and not leaving anything to be desired from the point of view of coherence and rigour, it must be admitted that it lacked in suppleness, and was hardly favourable for the development of numerical calculations and especially of algebraic calculations. What is more, its logical necessity could only be appreciated by minds taken with rigour and exercised in abstraction; it is therefore natural that with the decline of Greek mathematics, one sees reappearing little by little the ``naive'' point of view which had been kept by means of the tradition of the logisticians; it is that which dominates for example with Diophantus {[}91 a{]}, a true follower of this tradition rather than official Greek science; this latter, while reproducing for form's sake the Euclidean definition of number, understands really by the word ``number'', the unknown of the algebraic problems whose solution is, either an integer, or a fractional number, or even an irrational. \(^{7}\) Although this change in attitude, about the subject of number, was linked to one of the most important advances in the history of mathematics, namely the development of Algebra (see p.~48), it does not constitute, be it well understood, progress in itself, but rather regress.

It is not possible for us to follow here all the vicissitudes of the idea of number through Hindu, Arab and Occidental mathematics until the end of the Middle Ages: it is the ``naive'' notion of number that dominates; and, even though the Elements of Euclid were used as basis for the teaching of mathematics during this period, it is likely that the doctrine of Eudoxus remained generally not understood because the necessity for it did not arise any more. The ``ratios'' of Euclid were most often qualified as ``numbers''; the rules of calculation of integers were applied to them, obtaining thus exact results, without seeking to analyse deeply the reasons for the success of these methods.

We see however already R. Bombelli, in the middle of the XVIth century, writing on this subject, in his Algebra {[}28 b{]}, \(^{8}\) a point of view that (on condition that we assume acquired the results of book V of Euclid) is essentially correct; having recognised that once the unit of length is chosen there is a bijective correspondence between the lengths and the ratios of quantities, he defines the various algebraic operations (assuming the unit fixed, of course) on the lengths, and, representing numbers by lengths, obtains the geometric definition of the field of real numbers (a point of view for which the credit is most often given to Descartes), and gives thus to his Algebra a solid geometrical basis. \(^{9}\)

But the Algebra of Bombelli, even though singularly advanced for his times, did not go beyond the extraction of radicals and the solution by radicals of equations of the 2nd, 3rd and 4th degrees; of course the possibility of extraction of radicals is assumed by him without discussion. Simon Stévin {[}295{]}, also, adopts an analogous point of view on the subject of number, which is for him that which marks a measure of quantity, and which he considers as essentially ``continuous'' (without making precise the meaning he gives to this word); if he distinguishes between the ``geometrical numbers'' and the ``arithmetical numbers'', it is only from the accident of their mode of definition, without there being for him there a difference of kind; here besides is his last word on the subject: ``We conclude therefore that there are no absurd, irrational, irregular, inexplicable or deaf numbers; but that there is in them such excellence and agreement that we have subject matter for meditation night and day in their admirable perfection'' ({[}295{]}, p.~10). On the other hand, having been the first to have turned the tool of decimal fractions into a method of calculation, and proposed for them a notation already close to ours, he conceived clearly that these fractions furnish an indefinite approximation algorithm for every real number, as he brings out in his Appendice algebraique of 1594 ``containing a general rule for all Equations'' (a booklet of which the only known example was burnt at Louvain in 1914; but see {[}295{]}, p.~88). Such an equation having been put into the form \(P(x) = Q(x)\) (where \(P\) is a polynomial of degree greater than that of the polynomial \(Q\) , and \(P(0) < Q(0)\) ), we substitute for \(x\) the numbers 10, 100, 1000, \ldots{} until we find \(P(x) > Q(x)\) , which, he says, determines the number of digits of the root; then (if for example the root is found to have two digits) we substitute 10, 20, \ldots{} which determines the digit of tens; and then the same for the following digit, then for the successive decimal digits: ``And proceeding thus infinitely'', he says, ``one reaches infinitely near to that required'' ({[}295{]}, p.~88). As can be seen, Stévin had (no doubt the first) a clear idea of Bolzano's theorem and had recognised in this theorem the essential tool for the systematic solution of numerical equations; an intuitive conception of the numerical continuum is recognised there, at the same time, so clear that there remained little to do to make it definitively precise.

However, in the two centuries that followed, the definitive establishment of correct methods was twice slowed down by the development of two theories of which we do not have to give the history here: the infinitesimal calculus, and the theory of series. Across the discussions that they brought up, we recognise, as at all periods of the history of mathematics, the perpetual balancing between the seekers occupied with going forward, at the cost of some insecurity, persuaded that there will always be time later to consolidate the conquered territory, and the critical spirits, who (without necessarily giving an inch to the first as far as intuitive faculties and inventive talents go) do not believe they are wasting their efforts in devoting some effort to precise expression and to rigorous justification of their concepts. In the XVIIth century, the principal object of debate is the notion of the infinitely small, which, justified a posteriori by the results which it allows us to reach, seemed in open opposition to the axiom of Archimedes; and we see the most open spirits of the time finish by adopting a point of view little different from that of Bombelli, and which is distinguished from it mainly by the greater attention paid to the rigorous methods of the ancients; Isaac Barrow (the master of Newton, and who himself played an important part in the creation of the infinitesimal calculus) gives a brilliant exposé of it in his Lessons in Mathematics delivered at Cambridge in 1664-65-66 {[}16 a and b{]}; recognising the necessity, in order to find again for the subject of numbers the proverbial ``geometrical certainty'', of returning to the theory of Eudoxus, he presents at length, and very judicially, the defence of this latter (who, as he witnessed, appeared intelligible to many of his contemporaries) against those who tackled him with obscurity or even absurdity. On the other hand, defining numbers as symbols which denote the ratio of quantities, and susceptible of being combined amongst themselves by the operations of arithmetic, he obtains the field of real numbers, in terms taken up again after him by Newton in his Arithmetic and of which his successors until Dedekind and Cantor would not change anything.

But it is at about this time that the method of expansion in series was introduced, which soon, in the hands of impenitent algebraists, takes on an exclusively formal character and turns away the attention of mathematicians from the questions of convergence that the sane use of series raises in the domain of real numbers. Newton, the principal creator of this method, was still conscious of the necessity of considering these questions: and, if he had not sufficiently elucidated them, he had recognised at least that the power series that he was introducing converged ``most of the time'' at least as well as a geometric series (whose convergence was already known to the ancients) for small values of the variable ({[}233 a{]}, v. I, pp.~3-26); at about the same time, Leibniz had observed that an alternating series, with terms decreasing in absolute value and tending to 0, is convergent; in the following century, d'Alembert, in 1768, expresses doubts on the use of non-convergent series. But the authority of Bernoulli and especially Euler means that such doubts are exceptional at this time.

It is clear that the mathematicians who would have had the habit of using series in numerical calculation would never have neglected thus the notion of convergence; and it is not chance that the first who, in this domain as in many others, had brought about a return to correct methods, was a mathematician who, already from the prime of youth, had a love of numerical calculation: C. F. Gauss, who, almost a child, had used the algorithm of the arithmetico-geometric mean, \(^{10}\) could not avoid forming for himself a clear idea of limits; and we see him, in a fragment that dates from 1800 (but was only published in our time) ({[}124 a{]}, v X \(_1\) , p.~390), defining precisely, on the one hand the upper bound and the lower bound, on the other hand the upper limit and the lower limit of a sequence of real numbers; the existence of the first (for a bounded sequence) seeming to be assumed as obvious, and the latter having been correctly defined as limits for \(n\) tending to +∞, of \(\sup_{p≥0} u_{n+p}\) , \(\inf_{p≥0} u_{n+p}\) . Gauss, as well, gives also, in his memoir of 1812 on the hypergeometric series ({[}124 a{]}, v. III, p.~139) the first model of a discussion about convergence undertaken, as he says, ``in all rigour, and done to satisfy those whose preferences go to the rigorous methods of the ancient geometers'': it is true that this discussion, forming a secondary point in the memoir, does not go back to first principles in the theory of series; it is Cauchy who establishes these first, in his Cours d'Analyse of 1821 ({[}56 a{]}, (2), v. III), in a manner correct at all points, starting from the Cauchy criterion clearly stated, and assumed to be obvious; as, for the definition of number, he keeps to the point of view of Barrow and Newton, it can thus be said that for him the real numbers are defined by the axioms of quantities and the criterion of Cauchy: which indeed is sufficient to define them.

It is at the same time that another important aspect of the theory of real numbers is definitively clarified. As we have said, it has always been assumed as geometrically obvious that two continuous curves can not cross without intersecting; a principle that (made suitably precise) would be equivalent, it also, to the property of the straight line being a complete space. This principle is still part of the basis of the ``rigorous'' proof given by Gauss in 1799 of d'Alembert's theorem, according to which every polynomial with real coefficients admits a real or complex root ({[}124 a{]}, v. III, p.~1); the proof of the same theorem, given by Gauss in 1815 ({[}124 a{]}, v. III, p.~31), relies, in the same way as a previous attempt by Lagrange, on the principle, analogous but simpler, according to which a polynomial can not change sign without becoming zero (a principle that we have already seen used by Stévin). In 1817, Bolzano gives, starting with the criterion of Cauchy, a complete proof of this latter principle, which he obtains as a particular case of an analogous theorem about continuous numerical functions of a numerical variable {[}27 c{]}. Stating clearly (before Cauchy) the ``Cauchy criterion'', he seeks to justify it by an argument that, in the absence of all arithmetical definitions of a real number, was only and could only be a vicious circle; but, once this point is admitted, his work is entirely correct and quite remarkable, as containing, not only the modern definition of a continuous function (given here for the first time), with a proof of the continuity of polynomials, but even the proof of the existence of the lower bound of an arbitrary bounded set of real numbers (he does not speak of sets, but, which amounts to the same thing, of properties of real numbers). On the other hand, Cauchy, in his Cours d'Analyse ({[}56 a{]}, (2), v. III), defining, he also, continuous functions of one or more numerical variables, proves equally that a continuous function of one variable can not change sign without becoming zero, and this by the same argument as Simon Stévin, which naturally becomes correct, once continuity is defined, as soon as the Cauchy criterion is used (or as soon as the equivalent principle of ``nested intervals'' so-called is assumed, as Cauchy does at this point, of which the convergence of indefinite decimal fractions is of course only a particular case).

Once this point is reached, it only remained for mathematicians to make precise and develop the results acquired, by correcting a few errors and filling a few gaps. Cauchy, for example, had believed for a time that a convergent series, whose terms are continuous functions of a variable, has as sum a continuous function: the correction of this point by Abel, during the course of his important work on series ({[}1{]}, v. I, p.~219; cf.~also v. II, p.~257 and passim), ended finally by the elucidation by Weierstrass, in his lectures (unpublished, but which had a considerable influence), of the notion of uniform convergence (see p.~205). On the other hand, Cauchy had, without sufficient justification, assumed the existence of a minimum of a continuous function in one of the proofs given by him for the existence of roots of a polynomial; it is again

Weierstrass who brought clarity to questions of this kind by proving in his lectures the existence of a minimum for functions of numerical variables, defined on closed bounded intervals; it is following on from his critique of the unjustified application of this theorem to sets of functions (of which the ``Dirichlet principle'' is the best known example) that the movement of ideas begins which ends (see p.~143) with the general definition of compact spaces and the modern statement of the theorem.

At the same time, Weierstrass, in his lectures, had recognised the logical interest that there is in completely separating the idea of real numbers from the theory of quantities: to use this, in effect, reduces to defining axiomatically the set of points of the straight line (therefore definitively the set of real numbers) and assuming the existence of such a set; although this way of doing things is essentially correct, it is evidently preferable to start only with rational numbers, and to deduce the real numbers from there by completion. \(^{11}\) This is what was done, by different methods, and independently the one from the other, by Weierstrass, Dedekind, Méray and Cantor; whereas the procedure of ``cuts'' so-called, proposed by Dedekind ({[}79{]}, v. II, pp.~315-334) was much closer to the definitions of Eudoxus, the other proposed methods are nearer to the method used since Hausdorff to complete a metric space. It is at this moment also that Cantor starts to develop the theory of sets of real numbers, of which Dedekind had had the first idea {[}48{]}, obtaining thus the principal elementary results on the topology of the straight line, the structure of open sets, of closed sets, the notions of derived sets, perfect totally discontinuous sets, etc..

With Cantor, the theory of real numbers took, with few exceptions, its definitive form; let us indicate briefly in what way it has been extended. Apart from the work in general topology (see p.~143), and applications to Integration (see p.~221), it is mainly a case of research on the structure and classification of sets of points on the straight line and of numerical functions of real variables; they have their origin in the works of Borel {[}32 a{]}, aiming mainly at measure theory, but which ended up among other things at the definition of ``Borel sets'': these are sets belonging to the smallest family of subsets of R, containing the intervals, and closed with respect to union and countable intersection and the operation C. Intimately linked to these sets are the so-called ``Baire'' functions, that is to say those that can be obtained starting with continuous functions by the operation of limits, repeated ``transfinitely''; they were defined by Baire in the course of important work where he abandons entirely the measure point of view in order to tackle systematically the qualitative and ``topological'' aspect of these questions {[}11 a{]}: it is on this occasion that he defines and studies for the first time semi-continuous functions, and with a view to characterise limit functions of continuous functions, he introduces the important notion of a thin set (``first category'' set in the terminology of Baire) (see p.~165). As for the numerous works that have followed those of Baire, and which are due mainly to the Russian and especially Polish schools, we can only draw attention to their existence here (see p.~166).

\chapter{EXPONENTIALS AND LOGARITHMS.}

The history of the multiplicative group \(R_{+}^{*}\) of real numbers \textgreater0 is closely linked to that of the development of the notion of powers of a number \textgreater0, and of the notations employed to indicate them. The conception of a ``geometric progression'' made up of the successive powers of the same number goes back to the Egyptians and the Babylonians; it was familiar to Greek mathematicians, and one finds already in Euclid (Elements, Book IX, prop. 11) a general statement equivalent to the rule \(a^{m}a^{n}=a^{m+n}\) for integer exponents \textgreater0. In the Middle Ages, the French mathematician N. Oresme (XIVth century) finds this rule again; it is also with him that appears for the first time the notion of fractional exponent \textgreater0, with a notation already close to ours and some rules for calculation (stated in general form) concerning them (for example the two rules that we now write \((ab)^{1/n}=a^{1/n}b^{1/n},(a^{m})^{p/q}=(a^{mp})^{1/q}\) {[}74{]}. But the ideas of Oresme were too much in advance of the Mathematics of his time to exercise an influence on his contemporaries and his treatise sunk rapidly into forgetfulness. A century later, N. Chuquet stated again Euclid's rule; he introduces as well an exponential notation for the powers of unknowns in his equations, and does not hesitate to make use of the exponent 0 and integer exponents \textless0.¹ This time (and although the work of Chuquet remained in manuscript and does not seem to have been circulated much), the idea of isomorphism between the ``arithmetic progression'' of the exponents, and the ``geometric progression'' of the powers, will no longer be lost sight of; extended to negative exponents and to fractional exponents by Stifel ({[}298{]}, fol.~35 and 249-250), it ends up finally with the definition of logarithms and the construction of the first tables, undertaken independently by the Scot J. Napier, in 1614-1620, and the Swiss J. Bürgi (whose work did not appear until 1620, although his conception went back to the first years of the XVIIth century). With Bürgi, the continuity of the isomorphism established between R and \(R_{+}^{*}\) is implicitly assumed by the use of interpolation in the handling of the tables; it is on the contrary explicitly formulated in the definition of

Napier (as explicitly at least as was allowed by the fairly vague concept of continuity that obtained at that time). \(^{2}\)

We do not need to dwell here on the services rendered by logarithms in the numerical Calculus; from the theoretical point of view, their importance dates especially from the beginnings of the infinitesimal Calculus, with the discovery of series expansions of \(\log(1+x)\) and of \(e^{x}\) , and of the differential properties of these functions (see pp.~172 ff.). As far as the definition of exponentials and logarithms is concerned, it will be limited, until the middle of the XIXth century, to assuming intuitively the possibility of extending by continuity to the set of real numbers the function \(a^{x}\) defined for all rational x; and it is only once the notion of real number is definitively made precise and deduced from that of rational number, that there will be thought of giving a rigorous justification to this extension.

\chapter{N DIMENSIONAL SPACES.}

We have already had the opportunity of saying how the development of the analytic Geometry of the plane and space led mathematicians to introduce the notion of n dimensional space, which supplied them with a geometric language which was extremely convenient for expressing in a concise and simple way theorems of algebra concerning equations with an arbitrary number of variables, and notably all the general results of linear Algebra (see p.~61). But if, towards the middle of the XIXth century, this language had become current among numerous geometers, it remained purely conventional, and the absence of an ``intuitive'' representation of spaces of more than three dimensions seemed to forbid in these latter arguments ``by continuity'' that were allowed in the plane or in space based exclusively on ``intuition''. It is Riemann who was the first, in his research on Analysis situs and on the foundations of Geometry, who was bold to argue in that way by analogy with the case of three dimensional space (see p.~176); \(^{1}\) following him, numerous mathematicians set out to use, with great success, arguments of this nature, notably in the theory of algebraic functions of several complex variables. But the control of spatial intuition being very limited at that time, one could remain with good reason very sceptical as to the worth in proofs of such considerations, and only allow them purely heuristically, in as much as they made very plausible the exactness of certain theorems. It is thus that H. Poincaré, in his memoir of 1887 on the residues of double integrals of functions of two complex variables, avoids as far as possible all recourse to intuition in four dimensional space: ``As this hypergeometric language is repugnant still to many good minds'', he says, ``I will only make infrequent use of it''; the ``artifices'' that he employs to this end allow him to get back to topological arguments in three dimensional space, where he then hesitates no longer to appeal to intuition ({[}251 a{]}, v. III, pp.~443 ff.).

Elsewhere, the discoveries of Cantor, and notably the famous theorem establishing that R and \(R^{n}\) are equipotent (which seemed to put into question the very notion of dimension) \(^{2}\) showed that it was indispensable, in order to provide a solid basis for the arguments of Geometry and Topology, to free them entirely from all recourse to intuition. We have already said (cf.~p.~139) that this need is at the origin of the modern conception of general Topology; but even before the creation of this latter, a rigorous study had begun of the topology of number spaces and of their most immediate generalisations (the ``n dimensional varieties'') by methods arising mainly from the branch of Topology called ``combinatorial Topology'' or better ``algebraic Topology'', of which we can not speak here.
